% generated by roam.py handbook export -- не редактировать, править заметки
% 20261006125949-метопт_лекция_6_ускоренныи_метод_нестерова_и_квазиньютоновские_методы.org

\chapter{Ускоренный метод Нестерова и квазиньютоновские методы}\label{lecture:6}
\InputIfFileExists{text/06-intro}{}{}

\section{Ускоренный градиентный спуск}\label{lecture:6:1}  % конспект 5.1
\InputIfFileExists{text/06-1-intro}{}{}

\noteagain{4B497C50-C541-428D-94F6-B0F0DF464A50}{Оценивающая последовательность с квадратичной начальной функцией}  % -- лемма 5.3, без доказательства (повторение с лекции 5)

\noteagain{BB5E3FC2-069C-44F0-92D0-B4DC13092CCA}{Метод Нестерова}  % -- формулы из теоремы 5.1; форма без $v_k$ — рассуждение 5.1

\InputIfFileExists{text/C038105E-BD42-44A8-9542-2DBE9FFAAA3F-before}{}{}
% generated by roam.py handbook export -- не редактировать, править заметки
% 20261025144521-сходимость_ускоренного_метода_нестерова.org (body a7de0dadad65)

\begin{theorem}[{Сходимость ускоренного метода Нестерова}]\label[theorem]{note:C038105E-BD42-44A8-9542-2DBE9FFAAA3F}
  Пусть:

  \begin{itemize}
    \item $\left\{ \varphi_k \right\}$~--- \hyperref[note:4B497C50-C541-428D-94F6-B0F0DF464A50]{оценивающая последовательность с квадратичной начальной функцией},
    \item $\alpha_k = \sqrt{\frac{\mu}{L}}$
    \item $f$~--- $\mu$-\hyperref[note:22E546FC-186F-4A67-85AA-2C85DCEF9021]{сильно выпуклая функция}.
    \item $f \in C^{1,1}_L$~--- \hyperref[note:7DD1C7AB-1A49-4675-9F21-5A16C717F670]{Функции с липшицевым градиентом}
    \item $\varphi_0(x) = f(x_0) + \frac{\mu}{2}\|x - x_0\|^2$ (то есть $v_0 = x_0$, $\varphi_0^{*} = f(x_0)$), а точки $y_k$ в построении $\varphi_k$~--- те же, что ниже
  \end{itemize}

  Введем (\hyperref[note:BB5E3FC2-069C-44F0-92D0-B4DC13092CCA]{метод Нестерова})
  \[
  y_k = \frac{1}{1+\alpha_k} (\alpha_k v_k + x_k),\quad x_{k+1} = y_k - \frac{1}{L} \nabla f(y_k)
  \]

  Тогда $f(x_k) \le \varphi_k^{*}$
\end{theorem}

\begin{remark}
  За обозначениями $v_k, \varphi_k^{*}$ нужно идти \hyperref[note:4B497C50-C541-428D-94F6-B0F0DF464A50]{сюда}
\end{remark}

\begin{remark}
  При плохой обусловленности (то есть $\frac{\mu}{L}$ будет маленьким), мы получим более быструю сходимость, так как мы извлекаем из этого корень. Соответственно $\left(1 - \sqrt{\frac{\mu}{L}}\right)^k$ убывает быстрее, чем $\left(1 - \frac{\mu}{L}\right)^k$ у \hyperref[note:85D9D77D-13A2-42E0-BF93-AD085E0134B9]{градиентного спуска}: $\sqrt{\frac{\mu}{L}} > \frac{\mu}{L}$ при $\mu < L$.
\end{remark}

\paragraph*{Сходимость}
 Так как $f(x_k) \le \varphi^{*}_k$, то мы можем \hyperref[note:D21EBF43-7558-45EE-8198-60958E391E03]{оценить невязку через оценивающую последовательность}.

 Мы получим
 \[
f(x_k) - f^{*} \le \lambda_k(\varphi_0(x^{*}) - f^{*})
\]
Имея
\[
\lambda_k = (1-\alpha_k)\lambda_{k-1} = \left( 1 - \sqrt{\frac{\mu}{L}} \right) \lambda_{k-1} = \left( 1 - \sqrt{\frac{\mu}{L}} \right)^k
\]
Подставим его в неравенство
\[
f(x_k) - f^{*} \le \left( 1 - \sqrt{\frac{\mu}{L}} \right)^k  (\varphi_0(x^{*}) - f^{*})
\]

Получаем оценку, соответствующую \hyperref[note:DF2405BF-8B31-47EB-8DEF-B1A88034717C]{линейной сходимости} для \hyperref[note:C42D2C19-B35F-4C0E-AFEB-D880CBB72176]{невязки} $r_k = f(x_k) - f^{*}$

\begin{remark}
  Оценку можно записать без $\varphi$. При нашем $\varphi_0$
  \[
  \varphi_0(x^{*}) - f^{*} = f(x_0) - f^{*} + \frac{\mu}{2}\|x_0 - x^{*}\|^2 \le \frac{L + \mu}{2}\|x_0 - x^{*}\|^2,
  \]
  так как $f(x_0) - f^{*} \le \frac{L}{2}\|x_0 - x^{*}\|^2$ ($f \in C^{1,1}_L$, $\nabla f(x^{*}) = 0$). Итого
  \[
  f(x_k) - f^{*} \le \frac{L + \mu}{2}\left( 1 - \sqrt{\frac{\mu}{L}} \right)^k \|x_0 - x^{*}\|^2
  \]
\end{remark}

\begin{proof}
  Мы можем взять произвольную квадратичную функцию в качестве $\varphi_0$.
  Пусть $y_0 = x_0$ и $v_0 = x_0$
  \[
  \varphi_0 = f(x_0) + \frac{\mu}{2} \|x - v_0\|^2 = f(x_0) + \frac{\mu}{2} \|x - x_0\|^2
  \]
  Тогда $\varphi^{*}_0 = f(x_0)$

  Докажем неравенство по индукции.
  База $k=0$
  \[
  f(x_0) \le  \varphi^{*}_0 = f(x_0)
  \]
  Переход $k \mapsto k+1$ \\
  Из леммы о \hyperref[note:4B497C50-C541-428D-94F6-B0F0DF464A50]{квадратичной оценивающей последовательности}, мы знаем
  \begin{multline*}
  \varphi_{k+1}^{*} = \underline{(1-\alpha_{k}) \varphi^{*}_k} + \alpha_k f(y_k) - \frac{\alpha_k^2}{2\mu} \| \nabla f(y_k) \|^2 +\\
  \alpha_k (1-\alpha_k) \left( \frac{\mu}{2} \| y_k -v_k\|^2 + \left\langle \nabla f(y_k); v_k -y_k \right\rangle \right)
  \end{multline*}
  Заметим, что $\varphi^{*}_k \ge f(x_k)$ по индукционному предположению. Применим это к первому слагаемому.
  Получим
  \begin{multline*}
  \varphi_{k+1}^{*} \ge\underline{(1-\alpha_k)f(x_k)} + \alpha_k f(y_k) - \frac{\alpha_k^2}{2\mu} \| \nabla f(y_k) \|^2+\\
   \alpha_k (1-\alpha_k) \left( \underline{\frac{\mu}{2} \| y_k -v_k\|^2} + \left\langle \nabla f(y_k); v_k -y_k \right\rangle \right)
  \end{multline*}

  Также по \hyperref[note:1CA14F2C-637F-41DE-9E97-72075D19E97F]{выпуклости} $f$
  \[
  f(x_k) \ge f(y_k) + \left\langle \nabla f(y_k); x_k - y_k \right\rangle
  \]
  После использования неравенства давайте сгруппируем $(1-\alpha_k) f(y_k) + \alpha_k f(y_k) = f(y_k)$ и оценим снизу $\frac{\mu}{2} \|y_k - v_k\|^2 \ge 0$ \\
  Тогда получим
  \begin{multline*}
  \varphi_{k+1}^{*} \ge f(y_k) + (1-\alpha_k) \left\langle \nabla f(y_k); x_k- y_k \right\rangle - \frac{\alpha_k^2}{2\mu} \| \nabla f(y_k) \|^2+\\
   \alpha_k (1-\alpha_k) \left\langle \nabla f(y_k); v_k -y_k \right\rangle
  \end{multline*}
  Давайте по линейности совместим скалярные произведения, получим
  \[
  \varphi_{k+1}^{*} \ge f(y_k) + (1-\alpha_k) \left\langle \nabla f(y_k); x_k- y_k + \alpha_k v_k - \alpha_k y_k  \right\rangle - \frac{\alpha_k^2}{2\mu} \| \nabla f(y_k) \|^2
  \]

  Давайте посмотрим на  $\left\langle \nabla f(y_k); x_k- y_k + \alpha_k v_k - \alpha_k y_k  \right\rangle$ \\
  Мы задали
  \[
  y_k = \frac{1}{1+\alpha_k}(x_k + \alpha_k v_k)
  \]
  Заметим, что
  \[
  x_k - y_k + \alpha_k v_k - \alpha_ky_k  = x_k + \alpha_kv_k - (\alpha_k + 1)y_k =
  x_k + \alpha_k v_k - (x_k + \alpha_k v_k) = 0
  \]
  Тем самым скалярное произведение равно 0. Тогда
  \[
  \varphi_{k+1}^{*} \ge f(y_k) - \frac{\alpha_k^2}{2\mu} \| \nabla f(y_k) \|^2
  \]

  Мы знаем, что
  \[
  x_{k+1} = y_k - \frac{1}{L} \nabla f(y_k)
  \]
  Тогда по \hyperref[note:FFEA1F1D-B70C-4B4A-9386-8344A2A3BBC4]{оценке убывания за шаг градиентного спуска}
  \[
  f(x_{k+1}) \le f(y_k) - \frac{1}{L} \left( 1 - \frac{L}{2L} \right) \|\nabla f(y_k)\|^2 =
  f(y_k) - \frac{1}{2L} \|\nabla f(y_k)\|^2
  \]
  Получили оценку на $f(y_k)$
  \[
  f(y_k) \ge f(x_{k+1}) + \frac{1}{2L}\|\nabla f(y_k)\|^2
  \]
  Получили, что
  \[
  \varphi_{k+1}^{*} \ge f(x_{k+1}) + \frac{1}{2L} \|\nabla f(y_k)\|^2 - \frac{\alpha_k^2}{2\mu} \| \nabla f(y_k) \|^2
  \]
  Мы знаем, что $\alpha_k =\sqrt{\frac{\mu}{L}}$, тогда
  \[\frac{\alpha_k^2}{2\mu} = \frac{1}{2L}\]
  Мы получили, что нормы тоже сокращаются. Итого
  \[
  \varphi^{*}_{k+1} \ge f(x_{k+1})
  \]
  Индукция доказана.
\end{proof}
  % Сходимость ускоренного метода Нестерова -- теорема 5.1, следствие 5.1

\section{Метод тяжёлого шарика}\label{lecture:6:2}  % конспект 5.2
\InputIfFileExists{text/06-2-intro}{}{}

\InputIfFileExists{text/C6A43F85-13E0-46B3-96A8-BCB9F34ECD57-before}{}{}
% generated by roam.py handbook export -- не редактировать, править заметки
% 20261025183020-метод_тяжелого_шарика.org (body 98992cece584)

\begin{method}[{Метод тяжёлого шарика}]\label[method]{note:C6A43F85-13E0-46B3-96A8-BCB9F34ECD57}
  Метод с моментом (инерцией):
  \[
  x_{k+1} = x_k - \alpha \nabla f(x_k) + \beta (x_k - x_{k-1})
  \]
  где $\alpha > 0$~--- шаг, $\beta \in [0;1)$~--- коэффициент момента; нужны две начальные точки $x_0, x_1$. Слагаемое $\beta(x_k - x_{k-1})$~--- инерция шарика, катящегося по графику: он продолжает двигаться в прежнем направлении.
\end{method}

\begin{remark}
  В отличие от \hyperref[note:BB5E3FC2-069C-44F0-92D0-B4DC13092CCA]{метода Нестерова} в форме без $v_k$, градиент считается в самой точке $x_k$, а не в точке после инерционного сдвига.
\end{remark}

\paragraph*{Теорема Поляка (без доказательства)}
Пусть $f \in C^2$, $x^{*}$~--- невырожденная точка локального минимума: $\mu E \le \nabla^2 f(x^{*}) \le LE$, $\mu > 0$. Пусть $0 \le \beta < 1$, $0 < \alpha < \frac{2(1+\beta)}{L}$.

Тогда найдутся $q \in [0;1)$ и $\varepsilon > 0$ такие, что для любых $x_0, x_1$ с $\|x_0 - x^{*}\| < \varepsilon$, $\|x_1 - x^{*}\| < \varepsilon$ метод сходится к $x^{*}$ и
\[
\forall \delta \in (0; 1-q)\ \exists C(\delta):\ \|x_k - x^{*}\| \le C(\delta)(q + \delta)^k
\]
Наименьшее $q$ достигается при
\[
\alpha^{*} = \frac{4}{(\sqrt{L} + \sqrt{\mu})^2},\quad \beta^{*} = \left( \frac{\sqrt{L} - \sqrt{\mu}}{\sqrt{L} + \sqrt{\mu}} \right)^2,\quad q^{*} = \frac{\sqrt{L} - \sqrt{\mu}}{\sqrt{L} + \sqrt{\mu}}
\]

\begin{remark}
  $q^{*} = \frac{1 - \sqrt{\mu/L}}{1 + \sqrt{\mu/L}} \approx 1 - 2\sqrt{\frac{\mu}{L}}$~--- та же зависимость от $\sqrt{\frac{\mu}{L}}$, что в \hyperref[note:C038105E-BD42-44A8-9542-2DBE9FFAAA3F]{сходимости метода Нестерова}, и лучше, чем $1 - \frac{\mu}{L}$ у \hyperref[note:85D9D77D-13A2-42E0-BF93-AD085E0134B9]{градиентного спуска}. Но гарантия только локальная.
\end{remark}
  % Метод тяжёлого шарика -- рассуждение 5.2, теорема 5.2 (Поляк), без доказательства

\section{Квазиньютоновские методы}\label{lecture:6:3}  % конспект 7.2
\InputIfFileExists{text/06-3-intro}{}{}

\noteagain{C0ED5501-CFC6-42BA-A53A-1B346CBA545D}{Метод Ньютона}  % -- плюсы и минусы, рассуждение 7.1

\InputIfFileExists{text/6BA18878-DE34-46A5-ABE2-EDA09CAE56A7-before}{}{}
% generated by roam.py handbook export -- не редактировать, править заметки
% 20261025161604-квазиньютоновские_методы.org (body deb031ca6fc4)

\begin{method}[{Квазиньютоновские методы}]\label[method]{note:6BA18878-DE34-46A5-ABE2-EDA09CAE56A7}
  Основаны на \hyperref[note:C0ED5501-CFC6-42BA-A53A-1B346CBA545D]{методе Ньютона}.
  Они призваны сохранить его плюсы. Но при этом минимизировать его минусы, как обращение матрицы (которая затратна по скорости вычисления) и хранение \hyperref[note:592AD816-1D80-48E6-8399-8524FBC1E8F3]{гессиана} (много памяти нужно)
\end{method}

\paragraph*{Общая идея}
Тут мы используем последовательность матриц, а не просто матрицу вторых производных.
Мы как бы стараемся приблизить вторую производную этими матрицами.
\[
x_{k+1} = x_k - B_k^{-1} \nabla f(x_k),\quad f(x) \approx f(x_k) + \left\langle \nabla f(x_k); x -x_k \right\rangle + \frac{1}{2} \left\langle B_k(x-x_k); x-x_k \right\rangle
\]

Хотим, чтобы $B_{k+1} = F(B_k, ...) \Rightarrow B_{k+1}^{-1} = G (B_k^{-1}, ...)$

Какие условия накладываются на $B_k$?
\[
B_k = B_k^T\quad B_k > 0\quad y_k = B_{k+1} d_k
\]
где $y_k = B_{k+1}d_k$~--- \hyperref[note:E976DE49-0091-45E1-B621-A236F8B87859]{Условие секущей}

Рассмотрим $B_k$ как неизвестную. Ее задают $\frac{n(n+1)}{2}$ неизвестных скаляров.
Посчитаем сколько уравнений на эти параметры мы получим из наших условий:

\begin{itemize}
  \item Условие положительности дает $n$ неравенств по критерию Сильвестра
  \item Условие секущей дает тоже $n$ уравнений
\end{itemize}

Получается неполная система. Каждый метод сам определяет, как дополнить эту систему.
Отсюда и выходит разнообразие методов.

\paragraph*{Методы}
\begin{itemize}
  \item \hyperref[note:01306845-D37F-480A-8DCE-EE9CA022DA1A]{BFGS}
  \item \hyperref[note:F7506DC4-37FB-4E88-8178-B0FC678233C0]{SR1}
\end{itemize}
  % Квазиньютоновские методы -- рассуждение 7.1, замечание 7.4

\InputIfFileExists{text/E976DE49-0091-45E1-B621-A236F8B87859-before}{}{}
% generated by roam.py handbook export -- не редактировать, править заметки
% 20261025172143-условие_секущеи.org (body 537874588c71)

\begin{definition}[{Условие секущей}]\label[definition]{note:E976DE49-0091-45E1-B621-A236F8B87859}
  Пусть $f \in C^2$, $\left\{ x_k \right\}$~--- точки \hyperref[note:6BA18878-DE34-46A5-ABE2-EDA09CAE56A7]{квазиньютоновского метода} с матрицами $B_k$ и
  \[
  d_k:= x_{k+1} - x_k,\quad y_k = \nabla f(x_{k+1}) - \nabla f(x_k)
  \]

  Тогда для $B_{k+1}$ выполнено условие секущей, если
  \[
  y_k = B_{k+1}d_k
  \]

  Откуда оно берется:
  \[
  y_k = \nabla f(x_{k+1}) - \nabla f(x_k) = \int\limits_0^1 \nabla^2 f(x_k + \tau (x_{k+1} - x_k))\,d\tau(x_{k+1} - x_k) = \int\limits_0^1 \nabla^2 f(x_k + \tau d_k)\,d\tau d_k
  \]
  Получается, что
  \[
  y_k = \int\limits_0^1 \nabla^2f(x_k+\tau d_k)\, d\tau d_k = B_{k+1} d_k
  \]
  Получается, что $B_{k+1}$ должна оценивать приращение производной на направлении $d_k$ так же хорошо, как интеграл от $\nabla^2 f$ (можно наверное как-то более красиво сформулировать)
\end{definition}

\paragraph*{Необходимое условие: $d_k^T y_k > 0$}
Домножим исходное условие на $d_k^T$, получим
\[
d_k^Ty_k = d^T_kB_{k+1}d_k
\]
Так как $B_{k+1} > 0$, то при $d_k \ne 0$ получаем $d^T_kB_{k+1}d_k >0 \Rightarrow d^T_ky_k > 0$

Когда оно выполнено: \hyperref[note:DE298A0A-AF3E-4091-9587-64905F3E1DFB]{положительность скалярного произведения приращений аргумента и градиента}.
  % Условие секущей -- утверждение 7.1, замечание 7.3

\InputIfFileExists{text/DE298A0A-AF3E-4091-9587-64905F3E1DFB-before}{}{}
% generated by roam.py handbook export -- не редактировать, править заметки
% 20261025182915-положительность_скалярного_произведения_приращении_аргумента_и_градиента.org (body 9bc1ac5849a8)

\begin{theorem}[{Положительность скалярного произведения приращений аргумента и градиента}]\label[theorem]{note:DE298A0A-AF3E-4091-9587-64905F3E1DFB}
  Пусть $f \in C^1$,
  \[
  d_k = x_{k+1} - x_k,\quad y_k = \nabla f(x_{k+1}) - \nabla f(x_k)
  \]
  (как в \hyperref[note:E976DE49-0091-45E1-B621-A236F8B87859]{условии секущей}), $d_k \ne 0$. Если выполнено одно из условий:

  \begin{enumerate}
    \item $f$~--- $\mu$-\hyperref[note:22E546FC-186F-4A67-85AA-2C85DCEF9021]{сильно выпуклая};
    \item $x_{k+1} = x_k - \alpha_k B_k^{-1} \nabla f(x_k)$, где $B_k > 0$, а $\alpha_k$~--- шаг по \hyperref[note:5330B23D-95DD-4D5D-B4A3-63EBADB1DB7D]{условиям Армихо-Вульфа},
  \end{enumerate}

  то $d_k^Ty_k > 0$.
\end{theorem}

\begin{remark}
  $d_k^Ty_k > 0$~--- \hyperref[note:E976DE49-0091-45E1-B621-A236F8B87859]{необходимое условие} для условия секущей при $B_{k+1} > 0$. Здесь~--- когда оно выполнено.
\end{remark}

\begin{proof}
  \proofpart{Пункт 1}
  По \hyperref[note:1B5A3E25-56E8-4D5E-9CF6-E2731DA77604]{дифференциальному критерию сильной выпуклости} получаем
  \[
  f(x) \ge f(y) + \left\langle \nabla f(y); x-y \right\rangle + \frac{\mu}{2} \|x-y\|^2
  \]
  \[
  f(y) \ge f(x) + \left\langle \nabla f(x); y-x \right\rangle + \frac{\mu}{2} \|x-y\|^2
  \]
  Сложим эти два неравенства
  \[
  f(x) + f(y) \ge f(y) + f(x) + \left\langle \nabla f(y); x-y \right\rangle + \left\langle \nabla f(x); y-x \right\rangle + \mu \|x-y\|^2
  \]
  Сократим $f(x) + f(y)$ и заметим, что
  \begin{multline*}
  \left\langle \nabla f(y); x-y \right\rangle + \left\langle \nabla f(x); y-x \right\rangle = \\
   \left\langle \nabla f(y); x-y \right\rangle - \left\langle \nabla f(x); x-y \right\rangle =
  \left\langle \nabla f(y) - \nabla f(x); x-y \right\rangle
  \end{multline*}
  Тогда
  \[
  0 \ge \left\langle \nabla f(y) - \nabla f(x); x-y \right\rangle + \mu \|x-y\|^2
  \]
  Или
  \[
  \left\langle \nabla f(x) - \nabla f(y); x-y \right\rangle \ge \mu \|x-y\|^2
  \]

  Пусть $x = x_{k+1}, y=x_k$, тогда
  \[
  \left\langle \nabla f(x_{k+1}) - \nabla f(x_k); x_{k+1}-x_k \right\rangle \ge \mu \|d_k\|^2
  \]
  или
  \[
  \left\langle y_k; d_k \right\rangle \ge \mu \|d_k\|^2 > 0,\quad d_k \ne 0
  \]

  \proofpart{Пункт 2}
  Пусть
  \[
  p_k = - B_k^{-1}\nabla f(x_k)
  \]
  Тогда
  \[
  x_{k+1} = x_k + \alpha_k p_k
  \]
  $p_k$~--- направление спуска: $B_k > 0 \Rightarrow B_k^{-1} > 0$, поэтому при $\nabla f(x_k) \ne 0$
  \[
  \left\langle \nabla f(x_k); p_k \right\rangle = - \left\langle B_k^{-1} \nabla f(x_k); \nabla f(x_k) \right\rangle < 0
  \]
  Значит, мы делаем \hyperref[note:DF0CC2BA-9234-47D3-8988-7A1DAFC9B114]{спуск с шагом Армихо-Вульфа} по направлению $p_k$ (там оно обозначено $d_k$), и из условия Вульфа оттуда
  \[
  \left\langle \nabla f(x_{k+1}); p_k \right\rangle \ge  - c_2 \left| \left\langle \nabla f(x_k); p_k \right\rangle \right|
  \]
  Под модулем стоит $\left\langle \nabla f(x_k); p_k \right\rangle < 0$, раскроем модуль.
  \[
  \left\langle \nabla f(x_{k+1}); p_k \right\rangle \ge c_2 \left\langle \nabla f(x_k); p_k \right\rangle
  \]
  Вычтем $\left\langle \nabla f(x_k); p_k \right\rangle$ с обеих сторон
  \[
  \left\langle \nabla f(x_{k+1}) - \nabla f(x_k); p_k \right\rangle \ge (c_2-1) \left\langle \nabla f(x_k); p_k \right\rangle
  \]
  Заменим разность градиентов на $y_k$
  \[
  \left\langle y_k; p_k \right\rangle \ge (c_2-1) \left\langle \nabla f(x_k); p_k \right\rangle
  \]

  Заметим, что
  \[d_k = x_{k+1} - x_k = -\alpha_kB_k^{-1}\nabla f(x_k) = \alpha_kp_{k}\].
  Тогда
  \[d^T_ky_k = \left\langle y_k; d_k \right\rangle = \left\langle y_k; \alpha_kp_k \right\rangle\].
  Давайте домножим неравенство на $\alpha_k > 0$
  \[
  d_k^Ty_k = \alpha_k\left\langle y_k; p_k \right\rangle \ge  \alpha_k(c_2-1) \left\langle \nabla f(x_k); p_k \right\rangle > 0
  \]

  Оно больше нуля так как
  \[
  \alpha_k > 0\quad c_2-1 < 0\quad \left\langle \nabla f(x_k);p_k \right\rangle < 0
  \]
\end{proof}
  % Положительность скалярного произведения приращений аргумента и градиента -- утверждение 7.2 (в конспекте доказан п. 1, на лекции п. 2 — упражнение)

\InputIfFileExists{text/01306845-D37F-480A-8DCE-EE9CA022DA1A-before}{}{}
% generated by roam.py handbook export -- не редактировать, править заметки
% 20261025171513-bfgs.org (body 433253e9b0d7)

\begin{method}[{BFGS}]\label[method]{note:01306845-D37F-480A-8DCE-EE9CA022DA1A}
  Это \hyperref[note:6BA18878-DE34-46A5-ABE2-EDA09CAE56A7]{квазиньютоновский метод}.
  Задаем $B_k$, как
  \[
  B_k = J_kJ_k^T
  \]
  Заметим, что по построению $B_k$~--- симметрична и $\det J_k \ne 0 \Rightarrow B_k > 0$

  $J_{k+1}$ является решением задачи

  \[\left\{\begin{array}{l}
  \|J_{k+1} - J_k\| \rightarrow \min \\
  y_k = J_{k+1}J_{k+1}^Td_k
  \end{array}\right.
  \]

  Последнее~--- это \hyperref[note:E976DE49-0091-45E1-B621-A236F8B87859]{условие секущей}.
\end{method}
  % BFGS -- идея из опр. 8.1, без формулы пересчёта

\InputIfFileExists{text/F7506DC4-37FB-4E88-8178-B0FC678233C0-before}{}{}
% generated by roam.py handbook export -- не редактировать, править заметки
% 20261025171909-sr1.org (body f4e008a37faa)

\begin{method}[{SR1}]\label[method]{note:F7506DC4-37FB-4E88-8178-B0FC678233C0}
  \hyperref[note:6BA18878-DE34-46A5-ABE2-EDA09CAE56A7]{Квазиньютоновский метод} с матрицей по правилу
  \[B_{k+1} = B_k + \sigma_k v_k v_k^T\]
   где $\sigma_k \in \R$, $v_k \in \R^n$: поправка $\sigma_kv_kv_k^T$~--- симметричная матрица ранга 1 (отсюда название: Symmetric Rank 1), а $\sigma_k$, $v_k$ подбираются так, чтобы $B_{k+1}$ удовлетворяла \hyperref[note:E976DE49-0091-45E1-B621-A236F8B87859]{условию секущей}.
\end{method}

\paragraph*{Как подобрать $\sigma_k$ и $v_k$}
Подставим в условие секущей $y_k = B_{k+1}d_k$:
\[
\sigma_k v_k \left( v_k^T d_k \right) = y_k - B_k d_k
\]
Значит, $v_k \parallel y_k - B_k d_k$. Возьмём $v_k = y_k - B_k d_k$, тогда $\sigma_k = \frac{1}{\left\langle y_k - B_k d_k; d_k \right\rangle}$ и
\[
B_{k+1} = B_k + \frac{(y_k - B_k d_k)(y_k - B_k d_k)^T}{\left\langle y_k - B_k d_k; d_k \right\rangle}
\]

\begin{remark}
  Знаменатель может быть близок к нулю, а $B_{k+1} > 0$ не гарантируется~--- в отличие от \hyperref[note:01306845-D37F-480A-8DCE-EE9CA022DA1A]{BFGS}.
\end{remark}
  % SR1 -- нет в конспекте

\InputIfFileExists{text/06-outro}{}{}
